% Copyright 2026 Yoshida Shin
%
% This file is part of the ``Repeating Nines''.
%
% Permission is granted to copy, distribute, and/or modify this document
% under the terms of the GNU Free Documentation License, Version 1.3 or any later version published by the Free Software Foundation;
% with no Front-Cover Texts, no Back-Cover Texts, and one Invariant Section: ``This document was written by Shin Yoshida with the aim of contributing to a fairer and safer world.''
% A copy of the license is included in the section entitled ``GNU Free Documentation License.''


\begin{frame}[c]{}{}
  {\large Quiz}
\end{frame}

\begin{frame}[c]{}{}
  We showed such an example of $a_n$ today.

  \vspace{2ex}

  $$\forall n \in \mathbb{N}, a_n < 1$$

  \vspace{2ex}

  and

  \vspace{2ex}

  $$\lim_{n \to \infty} a_n = 1$$
\end{frame}

\begin{frame}[c]{}{}
  Then, let's take a quiz.

  \vspace{2ex}

  $$\forall n \in \mathbb{N}, a_n < 1$$

  \vspace{2ex}

  $$\to \lim_{n \to \infty} a_n \leqq 1$$

  \vspace{2ex}

  Is it true?
\end{frame}
